% GENERATED FILE. Edit canonical lessons, then run npm run generate:books.
% canonical-source-hash: fnv1a64:2185c7f898ba5638
% canonical-lessons: SA-C438-atapah, SA-C438-nidaghah, SA-C438-prakharah, SA-C438-svedah, SA-C438-dipyate

\chapter{Summer Heat}
\label{ch:sa-summer-heat}
\hlchaptermodality{438}

\begin{chapteropening}
\textbf{By the end of this chapter:} \emph{I can talk about the hot season and the blazing sun.}

Complete the last lesson of chapter 438: I can talk about the hot season and the blazing sun.
\end{chapteropening}

\section[\texorpdfstring{\sk{आतपः} (ātapaḥ)}{आतपः (ātapaḥ)}]{\sk{आतपः} (ātapaḥ) — sunshine; the sun's heat}

\label{lesson:SA-C438-atapah}



\begin{quote}
Before the new one: say the Sanskrit for a drop, then the Sanskrit for drips.
\end{quote}

\subsection*{You'll want to know: \sk{आतपः}}
\textbf{\sk{आतपः}} — \emph{ātapaḥ} — \textquotedblleft{}sunshine; the sun's heat\textquotedblright{}.

\begin{quote}
\sk{मध्याह्ने} \sk{आतपः} \sk{उष्णः} \sk{भवति।}
\end{quote}

\begin{quote}
\emph{madhyāhne ātapaḥ uṣṇaḥ bhavati.}
\end{quote}

\begin{quote}
At noon the sunshine is hot.
\end{quote}

A masculine noun.

\begin{cousinweb}
\textbf{\sk{आतपः}} (\emph{ātapaḥ}) is built from \emph{ā} and the root √tap, \textquotedblleft{}to heat, to burn\textquotedblright{}, which also gives \emph{tapati}, \textquotedblleft{}heats\textquotedblright{}; MW gives \emph{ātapa} as heat, especially of the sun, sunshine (MW).
\end{cousinweb}

\begin{grammarlens}[title={What you've built}]
One more word for getting around.
\end{grammarlens}

\subsection*{Your turn}
\begin{itemize}
\raggedright
  \item \emph{Say it:} \emph{ātapaḥ}
  \item \emph{Say it:} \emph{ātapaḥ}, once more
  \item \emph{Say it:} \emph{madhyāhne ātapaḥ uṣṇaḥ bhavati.} — At noon the sunshine is hot.
  \item \emph{Recall:} say the Sanskrit for a drop, then the Sanskrit for drips, then say \emph{ātapaḥ} again
\end{itemize}

\begin{tcolorbox}[breakable,colback=teal!4,colframe=teal!35!black,title={Before you move on}]
What does \sk{आतपः} mean? (\textquotedblleft{}Sunshine; the sun's heat\textquotedblright{}.) Say it once more.
\end{tcolorbox}
\section[\texorpdfstring{\sk{निदाघः} (nidāghaḥ)}{निदाघः (nidāghaḥ)}]{\sk{निदाघः} (nidāghaḥ) — the hot season}

\label{lesson:SA-C438-nidaghah}



\begin{quote}
Before the new one: say the Sanskrit for drips, then the Sanskrit for sunshine.
\end{quote}

\subsection*{You'll want to know: \sk{निदाघः}}
\textbf{\sk{निदाघः}} — \emph{nidāghaḥ} — \textquotedblleft{}the hot season\textquotedblright{}.

\begin{quote}
\sk{निदाघे} \sk{कूपाः} \sk{शुष्यन्ति।}
\end{quote}

\begin{quote}
\emph{nidāghe kūpāḥ śuṣyanti.}
\end{quote}

\begin{quote}
In the hot season the wells dry up.
\end{quote}

A masculine noun.

\begin{cousinweb}
\textbf{\sk{निदाघः}} (\emph{nidāghaḥ}) is built from \emph{ni}, \textquotedblleft{}down\textquotedblright{}, and the root √dah, \textquotedblleft{}to burn\textquotedblright{}, which also gives \emph{dahati}, \textquotedblleft{}burns\textquotedblright{}; MW gives \emph{nidāgha} as heat, warmth and the hot season (MW).
\end{cousinweb}

\begin{grammarlens}[title={What you've built}]
One more word for getting around.
\end{grammarlens}

\subsection*{Your turn}
\begin{itemize}
\raggedright
  \item \emph{Say it:} \emph{nidāghaḥ}
  \item \emph{Say it:} \emph{nidāghaḥ}, once more
  \item \emph{Say it:} \emph{nidāghe kūpāḥ śuṣyanti.} — In the hot season the wells dry up.
  \item \emph{Recall:} say the Sanskrit for drips, then the Sanskrit for sunshine, then say \emph{nidāghaḥ} again
\end{itemize}

\begin{tcolorbox}[breakable,colback=teal!4,colframe=teal!35!black,title={Before you move on}]
What does \sk{निदाघः} mean? (\textquotedblleft{}The hot season\textquotedblright{}.) Say it once more.
\end{tcolorbox}
\section[\texorpdfstring{\sk{प्रखरः} (prakharaḥ)}{प्रखरः (prakharaḥ)}]{\sk{प्रखरः} (prakharaḥ) — scorching; very hot}

\label{lesson:SA-C438-prakharah}



\begin{quote}
Before the new one: say the Sanskrit for sunshine, then the Sanskrit for the hot season.
\end{quote}

\subsection*{You'll want to know: \sk{प्रखरः}}
\textbf{\sk{प्रखरः}} — \emph{prakharaḥ} — \textquotedblleft{}scorching; very hot\textquotedblright{}.

\begin{quote}
\sk{निदाघे} \sk{आतपः} \sk{प्रखरः} \sk{भवति।}
\end{quote}

\begin{quote}
\emph{nidāghe ātapaḥ prakharaḥ bhavati.}
\end{quote}

\begin{quote}
In the hot season the sunshine is scorching.
\end{quote}

\begin{cousinweb}
\textbf{\sk{प्रखरः}} (\emph{prakharaḥ}) is built from \emph{pra}, \textquotedblleft{}very\textquotedblright{}, and \emph{khara}, \textquotedblleft{}hard, sharp\textquotedblright{}; MW gives \emph{prakhara} as very hard or rough and very hot (MW).
\end{cousinweb}

\begin{grammarlens}[title={What you've built}]
One more word for getting around.
\end{grammarlens}

\subsection*{Your turn}
\begin{itemize}
\raggedright
  \item \emph{Say it:} \emph{prakharaḥ}
  \item \emph{Say it:} \emph{prakharaḥ}, once more
  \item \emph{Say it:} \emph{nidāghe ātapaḥ prakharaḥ bhavati.} — In the hot season the sunshine is scorching.
  \item \emph{Recall:} say the Sanskrit for sunshine, then the Sanskrit for the hot season, then say \emph{prakharaḥ} again
\end{itemize}

\begin{tcolorbox}[breakable,colback=teal!4,colframe=teal!35!black,title={Before you move on}]
What does \sk{प्रखरः} mean? (\textquotedblleft{}Scorching; very hot\textquotedblright{}.) Say it once more.
\end{tcolorbox}
\section[\texorpdfstring{\sk{स्वेदः} (svedaḥ)}{स्वेदः (svedaḥ)}]{\sk{स्वेदः} (svedaḥ) — sweat}

\label{lesson:SA-C438-svedah}



\begin{quote}
Before the new one: say the Sanskrit for the hot season, then the Sanskrit for scorching.
\end{quote}

\subsection*{You'll want to know: \sk{स्वेदः}}
\textbf{\sk{स्वेदः}} — \emph{svedaḥ} — \textquotedblleft{}sweat\textquotedblright{}.

\begin{quote}
\sk{आतपे} \sk{कृषकस्य} \sk{मुखे} \sk{स्वेदः} \sk{अस्ति।}
\end{quote}

\begin{quote}
\emph{ātape kṛṣakasya mukhe svedaḥ asti.}
\end{quote}

\begin{quote}
In the sunshine there is sweat on the farmer's face.
\end{quote}

A masculine noun.

\begin{cousinweb}
\textbf{\sk{स्वेदः}} (\emph{svedaḥ}) comes from the root √svid, \textquotedblleft{}to sweat\textquotedblright{}; MW gives \emph{sveda} as sweating, sweat, and also warmth (MW).
\end{cousinweb}

\begin{grammarlens}[title={What you've built}]
One more word for getting around.
\end{grammarlens}

\subsection*{Your turn}
\begin{itemize}
\raggedright
  \item \emph{Say it:} \emph{svedaḥ}
  \item \emph{Say it:} \emph{svedaḥ}, once more
  \item \emph{Say it:} \emph{ātape kṛṣakasya mukhe svedaḥ asti.} — In the sunshine there is sweat on the farmer's face.
  \item \emph{Recall:} say the Sanskrit for the hot season, then the Sanskrit for scorching, then say \emph{svedaḥ} again
\end{itemize}

\begin{tcolorbox}[breakable,colback=teal!4,colframe=teal!35!black,title={Before you move on}]
What does \sk{स्वेदः} mean? (\textquotedblleft{}Sweat\textquotedblright{}.) Say it once more.
\end{tcolorbox}
\section[\texorpdfstring{\sk{दीप्यते} (dīpyate)}{दीप्यते (dīpyate)}]{\sk{दीप्यते} (dīpyate) — blazes; glows}

\label{lesson:SA-C438-dipyate}



\begin{quote}
Before the new one: say the Sanskrit for scorching, then the Sanskrit for sweat.
\end{quote}

\subsection*{You'll want to know: \sk{दीप्यते}}
\textbf{\sk{दीप्यते}} — \emph{dīpyate} — \textquotedblleft{}blazes; glows\textquotedblright{}.

\begin{quote}
\sk{मध्याह्ने} \sk{सूर्यः} \sk{दीप्यते।}
\end{quote}

\begin{quote}
\emph{madhyāhne sūryaḥ dīpyate.}
\end{quote}

\begin{quote}
At noon the sun blazes.
\end{quote}

Past (it …): \sk{सूर्यः} \sk{दीप्तवान्}, \emph{sūryaḥ dīptavān}.

\begin{cousinweb}
\textbf{\sk{दीप्यते}} (\emph{dīpyate}) comes from the root √dīp, \textquotedblleft{}to blaze, to shine\textquotedblright{}, which also gives \emph{dīpa}, \textquotedblleft{}a lamp\textquotedblright{}; MW gives \emph{dīpyate} as to blaze, flare, shine, with the participle \emph{dīpta} (MW).
\end{cousinweb}

\begin{grammarlens}[title={What you've built}]
One more word for getting around.
\end{grammarlens}

\subsection*{Your turn}
\begin{itemize}
\raggedright
  \item \emph{Say it:} \emph{dīpyate}
  \item \emph{Say it:} \emph{dīpyate}, once more
  \item \emph{Say it:} \emph{madhyāhne sūryaḥ dīpyate.} — At noon the sun blazes.
  \item \emph{Recall:} say the Sanskrit for scorching, then the Sanskrit for sweat, then say \emph{dīpyate} again
\end{itemize}

\begin{tcolorbox}[breakable,colback=teal!4,colframe=teal!35!black,title={Before you move on}]
What does \sk{दीप्यते} mean? (\textquotedblleft{}Blazes; glows\textquotedblright{}.) Say it once more.
\end{tcolorbox}
